\section{\IfLanguageName{dutch}{Database laag}{Database layer}}
\label{ch:architecture-db}

De databases bestaan uit een continu geüpdatete kopie van de productieomgeving, zodat onderzoekers steeds over actuele data beschikken zonder de operationele systemen te beïnvloeden.

De database-instanties zijn uitsluitend bedoeld voor het uitlezen van data; wijzigingen aan de data zijn niet toegestaan. Om de veiligheid en integriteit te waarborgen is er geen directe connectie met de productie-databases. In plaats daarvan wordt gewerkt met gescheiden databronnen binnen een gecontroleerde omgeving.

Voor toegang tot de data wordt gebruikgemaakt van een gebruiker met uitsluitend read-only rechten. Deze gebruiker heeft enkel SELECT-permissies op vooraf gedefinieerde views of schema’s en beschikt niet over write- of DDL-rechten. Daarnaast worden extra restricties toegepast, zoals statement timeouts en limieten op het aantal rijen per query, om misbruik en overbelasting te voorkomen.

Indien de belasting van de productieomgeving gevoelig is, kan er gebruik worden gemaakt van een read-replica. Deze replica dient als databron voor bijvoorbeeld Jupyter Notebook-omgevingen, zodat analytische workloads geen impact hebben op de performantie van de operationele systemen.

Verder kunnen er specifieke views of API-klare datasets worden voorzien, zodat gebruikers nooit rechtstreeks met onderliggende tabellen werken. Dit zorgt voor extra abstractie, verhoogde veiligheid en een stabielere datastructuur voor analyse.
